\section{Ideas on developments}
\begin{itemize}
\item $\vert \cH\vert$ is a measure of the complexity of our neural netwrok, i.e. proportional to the number of parameters: the more complex the neural network, the more functions it will be able to model, hence the {\it richer} our hypotheses class

\item the bound tells us that, for a fixed test size, as we increase complexity the gap between test error ($R$) and training error ($\hat R$) increases. this is aligend with trainig error going to 0 and test error increasing due to overfitting

\item it does not exlpain precisely an initial decrease in both errors, and a successive divergence

\item can we say that the increase in model complexity is exponential in the number of parameters added, hence $\log\vert\cH\vert$ gives a linear dependence on the number of parameters? 

\item anyways, the crucial point for our developments is that this bound only predicts an increase in the gap, and not the second descent in the test error. also, for a fixed data size you don't have a meaningful bound for $\vert \cH\vert$ large. 

\item also, uno spoiler che mi ha fatto Claude è che la cardinalità non è una buona misura di complessità per noi, perché GD esplora una porzione molto ridotta di tutto lo spazio. il concetto è quello di counting vs. covering.

\end{itemize}
